% Fragment for main.tex: uses its theorem/proof environments and booktabs.
% This replacement distinguishes the published coverage input from the
% independently replayable finite arithmetic certificate.
\section{Exact six- and seven-point proofs with an explicit coverage input}
\label{sec:six-seven-improved-proof}

Put $N_n=n(n-1)/2$ and
\[
 C_6=\frac{6^{15}}{5^{10}},\qquad
 C_7=\frac{7^{21}}{6^{15}}.
\]
The proof has two separate inputs. The published Faug\`ere--Svartz
computation supplies coverage of all physical stationary shapes.
The rational certificates supplied here cover and compare every algebraic
candidate in the retained charts. The latter verification does not
recompute the former decomposition.

\begin{lemma}[Reduction to normalized coefficient charts]
\label{lem:six-seven-chart-reduction}
A global maximizer of $\Delta_n$ on $\sum_i|z_i|^2=n$ exists, is
collision-free, and has $\sum_i z_i=0$. If $w_i=\alpha z_i$ for
$\alpha\ne0$, and $P(t)=\prod_i(t-w_i)$, $Q=\sum_i|w_i|^2$, then
\begin{equation}\label{eq:six-seven-uniform-residual}
 QP''(w_i)-2N_n\overline{w_i}P'(w_i)=0\quad(1\le i\le n),
\end{equation}
and its value at moment $n$ is
\begin{equation}\label{eq:six-seven-uniform-value}
 \mathcal V_n(w)=\left(\frac nQ\right)^{N_n}
                    \prod_{i<j}|w_i-w_j|^2.
\end{equation}
Every such shape can be put in a chart in which the first nonzero
elementary symmetric coefficient is $1$; complex scaling loses no
physical shape and preserves \eqref{eq:six-seven-uniform-value}.
\end{lemma}
\begin{proof}
The constraint sphere is compact and the objective is continuous.
A regular $n$-gon gives a positive value, whereas a collision gives
zero, so a maximizer has distinct points. Let $c=n^{-1}\sum_i z_i$.
Translation by $-c$ preserves all pair differences and changes the
moment to $n-n|c|^2$. If $c\ne0$, rescaling the translated points to
moment $n$ strictly increases the positive objective, a contradiction.

The Lagrange equations for $\log\Delta_n$ at a maximizer are
\[
 F_i(z):=\sum_{j\ne i}\frac1{z_i-z_j}=\lambda\overline{z_i}.
\]
Pairwise summation gives
$\sum_i z_iF_i(z)=\sum_{i<j}(z_i-z_j)/(z_i-z_j)=N_n$.
Consequently $\lambda=N_n/n>0$. Under $w_i=\alpha z_i$,
$F_i(w)=\lambda|\alpha|^{-2}\overline{w_i}$, and the positive
force constant $\kappa=\lambda|\alpha|^{-2}$ obeys $\kappa Q=N_n$.
The identity $F_i(w)=P''(w_i)/(2P'(w_i))$ proves the residual equation.
Scaling a shape of moment $Q$ to moment $n$ multiplies every squared
pair difference by $n/Q$, proving the value formula.

The centered root polynomial is
$P(t)=t^n+\sum_{k=2}^n(-1)^ke_k t^{n-k}$.
Some $e_k$ is nonzero, since otherwise every root would be zero.
Choosing $\alpha$ with $\alpha^ke_k=1$ normalizes its first nonzero
coefficient, with all earlier coefficients still zero.
\end{proof}

\paragraph{The exact external coverage statement.}
We use the collision-saturated coefficient-chart computation reported in
Faug\`ere--Svartz~\cite[Sections~5.3--5.7]{FS}, with the archived
author bases~\cite{FSdata,FSdataSeven}, in the following precise form:
every collision-free physical stationary shape for $n=6$ or $n=7$,
up to complex scaling and permutation, is represented in the retained
charts specified below. This is a necessary coverage statement; the
charts may also contain nonphysical algebraic candidates. Theorem~5
of~\cite{FS} gives the symmetric equations, whereas the chart
classification is the subsequent published computation.
This coverage concerns solutions of the original physical equilibrium
equations, including unstable and collinear configurations; it is not a
classification restricted to stable equilibria. Sections~5.3--5.4 and
5.7 occur on proceedings pages~176 and~177, respectively.

For $n=6$ the author file supplies eleven components, of parameter
degrees
\[
 22,28,12,4,2,1,1,6,2,1,1,
\]
with total $80$. Its header $\texttt{N=7}$ is a metadata typo: the
linked file is $\texttt{vortices\_N6.txt}$ and its variables are
$e_2,\ldots,e_6$. Exact substitution into the six-point coefficient
equations supplies an additional integrity check. It proves forward
containment, and does not by itself prove completeness of these charts.

For $n=7$, Section~7 and Figure~4 of~\cite[p.~178]{FS} are also used. Their
$e_2=1$ component has degree $526$, decomposed into degrees
$6,18,356,82,64$. Figure~4 marks the degree-$356$ branch physically
empty. Its polynomial and its independent real-feasibility certificate
are absent from our supplied files. The retained file has seven charts,
of parameter degrees
\[
 1,1,2,6,18,82,64,
\]
with total $174$. The degree-$18$ branch, also marked physically
empty in the paper, is retained and independently excluded below.
Thus our seven-point verification includes all $174$ retained
coefficient candidates, and does not claim to include the missing
$356$ complex candidates. Their physical exclusion remains a cited
external input.

\begin{proposition}[Exact comparison on every retained chart]
\label{prop:six-seven-retained-comparison}
For every retained six-point candidate that satisfies
\eqref{eq:six-seven-uniform-residual},
$\mathcal V_6\le C_6$. Equality can occur only for
$P(t)=t(t^5-1)$.
For every retained seven-point candidate that satisfies the same
residual with $n=7$, $\mathcal V_7\le C_7$. Equality can occur only
for $P(t)=t(t^6+1)$.
\end{proposition}
\begin{proof}
The coefficient conventions are
\[
 \begin{split}
 P_6(t)&=t^6+e_2t^4-e_3t^3+e_4t^2-e_5t+e_6,\\
 P_7(t)&=t^7+e_2t^5-e_3t^4+e_4t^3-e_5t^2+e_6t-e_7.
 \end{split}
\]
The standard-library programs
$\texttt{certify\_fs6\_disks.py}$ and
$\texttt{certify\_fs7\_disks.py}$ parse the exact author polynomials
without executing them. They isolate every parameter root and every
root of the associated point polynomial in disjoint Gaussian-rational
disks. Each disk contains one root by a strict rational Rouch\'e
inequality, and the number of disks equals the polynomial degree.
Thus no parameter root or point root in any retained chart is omitted.
Coefficient uncertainty from the parameter disk is included throughout.

Here are the replayed six-point results. ``Bounded'' means a strict
upper bound on the normalized objective; it does not assert that the
candidate satisfies the physical equations.
\begin{center}
\begin{tabular}{rrrrl}
\toprule
Component&Parameters&Force excluded&Bounded&Upper endpoint\\
\midrule
1&22&20&2&$3170$\\
2&28&28&0&---\\
3&12&12&0&---\\
4&4&2&2&$42454$\\
5&2&2&0&---\\
6&1&0&1&$18699$\\
7&1&0&1&$18699$\\
8&6&4&2&$22400$, $133$\\
9&2&1&1&$7225$\\
10&1&0&0&$C_6$, attaining\\
11&1&0&1&$6^6$, exact\\
\bottomrule
\end{tabular}
\end{center}
There are $80$ parameter roots and $480$ point-root disks.
The $69$ force exclusions and $10$ strictly smaller candidates leave
one attaining chart. Every displayed upper endpoint is at most $6^6$,
and
\[
 C_6>6^6\quad\Longleftrightarrow\quad
 6^9=10077696>9765625=5^{10}.
\]

The seven-point results are
\begin{center}
\begin{tabular}{rrrrl}
\toprule
Chart&Parameters&Force excluded&Bounded&Upper endpoint\\
\midrule
1&1&0&1&$7^7=823543$, exact\\
2&1&0&0&$C_7$, attaining\\
3&2&1&1&$40253$\\
4&6&4&2&$159666$\\
5&18&18&0&---\\
6&82&78&4&$268962$\\
7&64&54&10&$822629$\\
\bottomrule
\end{tabular}
\end{center}
There are $174$ parameter roots and $1218$ point-root disks.
The $155$ force exclusions and $18$ strictly smaller candidates leave
one attaining chart. Every displayed upper endpoint is at most $7^7$,
and
\[
 C_7>7^7\quad\Longleftrightarrow\quad
 7^{14}=678223072849>470184984576=6^{15}.
\]
The force exclusions and objective bounds are justified uniformly on
the exact root disks by the estimates given next. All deciding checks
use explicit exceptions and remain active under Python $\texttt{-O}$.
\end{proof}

\paragraph{Root coverage, physical exclusions and value bounds.}
Write $\|c\|_1=|\Re c|+|\Im c|$ and
$\|c\|_\infty=\max\{|\Re c|,|\Im c|\}$ for Gaussian rationals.
Let $p_0$ be a nominal polynomial, let its coefficient of $t^j$ have
uncertainty at most $\epsilon_j$, and put
$a_k=p_0^{(k)}(c)/k!$. The test
\[
 \|a_0\|_1+\sum_{k\ge2}\|a_k\|_1r^k
       +\sum_j\epsilon_j(\|c\|_1+r)^j
 <r\|a_1\|_\infty
\]
compares the allowed polynomial on $|t-c|=r$ with its nominal
linear term $a_1(t-c)$. Rouch\'e's theorem therefore gives exactly
one root, counted with multiplicity. Pairwise disk disjointness and
the degree count prove exhaustive coverage in each retained chart.

Let $c_i$ be the point-root centers, of common radius $r$, and set
\[
 Q_0=\sum_i|c_i|^2,\qquad
 E_Q=\sum_i(2\|c_i\|_1r+r^2).
\]
Then $|Q-Q_0|\le E_Q$. Let $v_j=p_0^{(j)}(c_i)$ and let $E_j$
bound $|P^{(j)}(w_i)-v_j|$, including coefficient and root variation.
The actual physical residual differs from
$Q_0v_2-2N_n\overline{c_i}v_1$ by at most
\[
 E_Q\|v_2\|_1+(Q_0+E_Q)E_2
 +2N_n\{r\|v_1\|_1+(\|c_i\|_1+r)E_1\}.
\]
If the nominal residual's $\|\cdot\|_\infty$ exceeds this bound
at even one point, the entire coefficient candidate is nonphysical.
This tests the actual pointwise conjugate pairing, rather than only
equality of two unordered sets of roots.

For each surviving candidate, $Q_0>E_Q$ and
\begin{equation}\label{eq:six-seven-rational-upper}
 U_n=\left(\frac n{Q_0-E_Q}\right)^{N_n}
 \prod_{i<j}\bigl(|c_i-c_j|^2+4r\|c_i-c_j\|_1+4r^2\bigr)
\end{equation}
is a rational upper bound on \eqref{eq:six-seven-uniform-value}.
The tables give upward-rounded integer endpoints of these rational
bounds, or the exact regular-polygon values. Hence their strict
comparison with $C_n$ proves the proposition without needing to
establish physical sufficiency for any bounded candidate.

\begin{theorem}[Global values using the published physical coverage]
\label{thm:six-seven-global-improved}
For $\sum_i|z_i|^2=n$,
\[
 M_6=\frac{6^{15}}{5^{10}},\qquad
 M_7=\frac{7^{21}}{6^{15}}.
\]
Equality for $n=6$ holds exactly for the origin and a regular
pentagon of radius $\sqrt{6/5}$; equality for $n=7$ holds exactly
for the origin and a regular hexagon of radius $\sqrt{7/6}$,
in each case up to rotation and permutation.
\end{theorem}
\begin{proof}
Every global maximizer is a collision-free centered physical stationary
shape by Lemma~\ref{lem:six-seven-chart-reduction}. The stated external
coverage places it in a retained chart, and
Proposition~\ref{prop:six-seven-retained-comparison} gives the upper
bound with a strict gap for every chart other than the attaining one.

To verify attainment independently of the classification, put
$m=n-1$ and consider the roots of $P(t)=t(t^m-a)$, $|a|=1$.
They are the origin and a rotated unit regular $m$-gon, with moment
$m$ and squared pair-distance product $m^m$. At every nonzero root,
$P''(w)/(2P'(w))=n/(2w)$, and the force at the origin is zero.
Thus the residual equation holds directly. Scaling to moment $n$
gives
\[
 (n/m)^{N_n}m^m
 =\frac{n^{N_n}}{(n-1)^{(n-1)(n-2)/2}}.
\]
For $n=6,7$ this is respectively $C_6,C_7$.
The strict comparisons already prove that an attaining shape must
belong to this chart. Moment normalization fixes its radius and
the remaining complex phase is exactly an overall rotation, proving
both equality classifications.
\end{proof}

\paragraph{Reproducibility boundary.}
The exact replay counts are $80=69+10+1$ and $174=155+18+1$.
The six-point component-identity checker has also been changed from
$\texttt{assert}$ statements to explicit arithmetic checks, so its
decisions remain active under $\texttt{python -S -O}$.
An independent reconstruction of the saturated Gr\"obner decompositions
is not included, and the missing seven-point degree-$356$ physical
exclusion is not claimed to have been reproduced. A proof using the
published classification and an entirely standalone computational
classification are different reproducibility claims; the theorem here
uses the former.
